\documentclass[12pt]{article}
\usepackage{amsmath,amssymb}

\begin{document}
\section*{{2023 AIME II Problems/Problem 12}}
Source: AoPS Wiki

\subsection*{Solution}
Because $M$ is the midpoint of $BC$ , following from the Stewart's theorem, $AM = 2 \sqrt{37}$ . Because $A$ , $B$ , $C$ , and $P$ are concyclic, $\angle BPA = \angle C$ , $\angle CPA = \angle B$ . Denote $\theta = \angle PBQ$ . In $\triangle BPQ$ , following from the law of sines, \[ \frac{BQ}{\sin \angle BPA} = \frac{PQ}{\sin \angle PBQ} \] Thus, \[ \frac{BQ}{\sin C} = \frac{PQ}{\sin \theta} . \hspace{1cm} (1) \] In $\triangle CPQ$ , following from the law of sines, \[ \frac{CQ}{\sin \angle CPA} = \frac{PQ}{\sin \angle PCQ} \] Thus, \[ \frac{CQ}{\sin B} = \frac{PQ}{\sin \theta} . \hspace{1cm} (2) \] Taking $\frac{(1)}{(2)}$ , we get \[ \frac{BQ}{\sin C} = \frac{CQ}{\sin B} \] In $\triangle ABC$ , following from the law of sines, \[ \frac{AB}{\sin C} = \frac{AC}{\sin B} . \hspace{1cm} (3) \] Thus, Equations (2) and (3) imply \begin{align*} \frac{BQ}{CQ} & = \frac{AB}{AC} \\ & = \frac{13}{15} . \hspace{1cm} (4) \end{align*} Next, we compute $BQ$ and $CQ$ . We have \begin{align*} BQ^2 & = AB^2 + AQ^2 - 2 AB\cdot AQ \cos \angle BAQ \\ & = AB^2 + AQ^2 - 2 AB\cdot AQ \cos \angle BAM \\ & = AB^2 + AQ^2 - 2 AB\cdot AQ \cdot \frac{AB^2 + AM^2 - BM^2}{2 AB \cdot AM} \\ & = AB^2 + AQ^2 - AQ \cdot \frac{AB^2 + AM^2 - BM^2}{AM} \\ & = 169 + AQ^2 - \frac{268}{2 \sqrt{37}} AQ . \hspace{1cm} (5) \end{align*} We have \begin{align*} CQ^2 & = AC^2 + AQ^2 - 2 AC\cdot AQ \cos \angle CAQ \\ & = AC^2 + AQ^2 - 2 AC\cdot AQ \cos \angle CAM \\ & = AC^2 + AQ^2 - 2 AC\cdot AQ \cdot \frac{AC^2 + AM^2 - CM^2}{2 AC \cdot AM} \\ & = AC^2 + AQ^2 - AQ \cdot \frac{AC^2 + AM^2 - CM^2}{AM} \\ & = 225 + AQ^2 - \frac{324}{2 \sqrt{37}} AQ . \hspace{1cm} (6) \end{align*} Taking (5) and (6) into (4), we get $AQ = \frac{99}{\sqrt{148}}$ Therefore, the answer is $99 + 148 = \boxed{\textbf{(247) }}$ ~Steven Chen (Professor Chen Education Palace, www.professorchenedu.com)

\end{document}
